\subsection{模的正向极限}

\BourbakiRunInHeading{以下假设I是右有向的。}

设 \((\mathbf{A}_{\alpha}, \phi_{\beta\alpha})\) 是环的一个正向系统（I，§10，第3号）， \((\mathbf{E}_{\alpha}, f_{\beta\alpha})\) 是交换群（加法记法）的一个正向系统（I，§10，第3号），并假设每个 \(E_{\alpha}\) 具有一个左 \(A_{\alpha}\)-模结构；进而，假设对 \(\alpha \leqslant \beta\) , \((f_{\beta\alpha}, \phi_{\beta\alpha})\) 是 \(E_{\alpha}\) 到 \(E_{\beta}\) 的一个双态射（§1，第13号），换句话说，

\[
f _ {\beta \alpha} (\lambda_ {\alpha} x _ {\alpha}) = \phi_ {\beta \alpha} (\lambda_ {\alpha}) f _ {\beta \alpha} (x _ {\alpha})\tag{4}
\]

对 \(x_{\alpha} \in \mathrm{E}_{\alpha}, \lambda_{\alpha} \in \mathrm{A}_{\alpha}\)；则 \(\mathbf{E} = \varinjlim \mathbf{E}_{\alpha}\) 具有一个关于 \(\mathbf{A} = \varinjlim \mathbf{A}_{\alpha}\) 的左模结构（I，§10，第4号）。对所有 \(\alpha \in \mathbf{I}\)，设 \(f_{\alpha}: \mathrm{E}_{\alpha} \to \mathrm{E}\) , \(\phi_{\alpha}: \mathrm{A}_{\alpha} \to \mathrm{A}\) 为典范映射；则 \((f_{\alpha}, \phi_{\alpha})\) 是 \(\mathbf{E}_{\alpha}\) 到 \(\mathbf{E}\) 的一个双态射。我们将称 \((\mathbf{E}_{\alpha}, f_{\beta \alpha})\) 是左 \(\mathbf{A}_{\alpha}\)-模的一个正向系统，且A-模 \(\mathbf{E}\) 是其正向极限。

设 \((\mathbf{E}_{\alpha}^{\prime}, f_{\beta \alpha}^{\prime})\) 是左 \(A_{\alpha}\)-模的另一个正向系统，并且对所有 \(\alpha\)，设 \(u_{\alpha}: \mathrm{E}_{\alpha}' \to \mathrm{E}_{\alpha}\) 为 \(\mathbf{A}_{\alpha}\)-线性映射，且这些映射构成一个正向系统；则 \(u = \varinjlim u_{\alpha}\) 是从 \(\varinjlim \mathrm{E}_{\alpha}'\) 到 \(\varinjlim \mathrm{E}_{\alpha}\) 的一个A-线性映射。

此外：

命题3. 设 \((\mathbf{E}_{\alpha}, f_{\beta \alpha})\) , \((\mathbf{E}_{\alpha}', f_{\beta \alpha}')\) , \((\mathbf{E}_{\alpha}'', f_{\beta \alpha}'')\) 是 \(\mathbf{A}_{\alpha}\)-模的三个正向系统，且 \((u_{\alpha})\) , \((v_{\alpha})\) 是 \(\mathbf{A}_{\alpha}\)-线性映射的两个正向系统，使得序列

\[
\mathrm{E} _ {\alpha} ^ {\prime} \xrightarrow {u _ {\alpha}} \mathrm{E} _ {\alpha} \xrightarrow {v _ {\alpha}} \mathrm{E} _ {\alpha} ^ {\prime \prime}
\]

对所有 \(\alpha\) 正合。于是，记 \(u = \varinjlim u_{\alpha}\) , \(v = \varinjlim v_{\alpha}\)，则序列

\[
\lim _ {\longrightarrow} \mathrm{E} _ {\alpha} ^ {\prime} \xrightarrow {u} \lim _ {\longrightarrow} \mathrm{E} _ {\alpha} \xrightarrow {v} \lim _ {\longrightarrow} \mathrm{E} _ {\alpha} ^ {\prime \prime}
\]

是正合的。

\(u(\varinjlim \mathrm{E}_{\alpha}^{\prime}) = \varinjlim u_{\alpha}(\mathrm{E}_{\alpha}^{\prime})\) 和 \(v^{-1}(0) = \varinjlim v_{\alpha}^{-1}(0)\) （集合论，III，§ 7，第 6 号，命题 7 的推论）。

粗略地说，命题3也可表述为：取正向极限保持正合性。

命题4. 设 \((\mathbf{E}_{\alpha}, f_{\beta \alpha})\) 是 \(\mathbf{A}_{\alpha}\)-模的一个正向系统，\(\mathbf{E} = \varinjlim \mathbf{E}_{\alpha}\) 是其正向极限，\(\phi_{\alpha}: \mathbf{A}_{\alpha} \to \mathbf{A}\) 与 \(f_{\alpha}: \mathbf{E}_{\alpha} \to \mathbf{E}\) 为对所有 \(\alpha \in \mathbf{I}\) 的典范映射。若对所有 \(\alpha \in \mathbf{I}\)，\(\mathbf{S}_{\alpha}\) 是 \(\mathbf{E}_{\alpha}\) 的一个生成系，则 \(\mathbf{S} = \bigcup_{\alpha \in \mathbf{I}} f_{\alpha}(\mathbf{S}_{\alpha})\) 是 \(\mathbf{E}\) 的一个生成系。

每个 \(x \in \mathbf{E}\) 具有 \(f_{\alpha}(x_{\alpha})\) 的形式（对某个 \(\alpha \in \mathbf{I}\) 及某个 \(x_{\alpha} \in \mathbf{E}_{\alpha}\)），且由假设 \(x = \sum_{i} \lambda_{\alpha}^{(i)} y_{\alpha}^{(i)}\)，其中 \(\lambda_{\alpha}^{(i)} \in \mathbf{A}_{\alpha}\) 且 \(y_{\alpha}^{(i)} \in \mathbf{S}_{\alpha}\)；记 \(\lambda^{(i)} = \phi_{\alpha}(\lambda_{\alpha}^{(i)})\) , \(y^{(i)} = f_{\alpha}(y_{\alpha}^{(i)})\)，我们得到 \(x = \sum_{i} \lambda^{(i)} y^{(i)}\)。

命题5. 在命题4的假设与记号下，设对所有 \(\alpha \in I\)，\(E_{\alpha}\) 是子模族 \((M_{\alpha}^{\lambda})_{\alpha \in L}\) 的直和（指标集 \(L\) 与 \(\alpha\) 无关），且 \(f_{\beta \alpha}(M_{\alpha}^{\lambda}) \subset M_{\alpha}^{\lambda}\)（对 \(\alpha \leqslant \beta\) 及所有 \(\lambda \in L\)）。则 \(E\) 是子模 \(M^{\lambda} = \varinjlim_{\alpha} M_{\alpha}^{\lambda} (\lambda \in L)\) 的直和。

由命题4可知，E是诸 \(\mathbf{M}^{\lambda}\) 之和。设 \((y^{\lambda})_{\lambda \in L}\) 是一个族，使得 \(y^{\lambda} \in \mathbf{M}^{\lambda}\)（对所有 \(\lambda \in L\)）且其支撑有限，并设 \(\sum_{\lambda} y^{\lambda} = 0\)。由集合论，III，§7，第5号，引理1，存在一个 \(\alpha \in I\) 和一个有限支撑族 \((x_{\alpha}^{\lambda})_{\alpha \in L}\)（由 \(\mathbf{E}_{\alpha}\) 的元素组成），使得 \(x_{\alpha}^{\lambda} \in \mathbf{M}_{\alpha}^{\lambda}\) 且 \(y^{\lambda} = f_{\alpha}(x_{\alpha}^{\lambda})\)（对所有 \(\lambda \in L\)）。关系 \(f_{\alpha}\left(\sum_{\lambda \in L} x_{\alpha}^{\lambda}\right) = 0\) 蕴含存在一个 \(\beta \geqslant \alpha\)，使得 \(f_{\beta \alpha}\left(\sum_{\lambda \in L} x_{\alpha}^{\lambda}\right) = 0\)（集合论，III，§ 7，第 5 号，引理 1），这可写成 \(\sum_{\lambda \in L} x_{\beta}^{\lambda} = 0\)，其中由假设 \(x_{\beta}^{\lambda} = f_{\beta \alpha}(x_{\alpha}^{\lambda}) \in \mathrm{M}_{\beta}^{\lambda}\)；因此 \(x_{\beta}^{\lambda} = 0\)（对所有 \(\lambda \in \mathbf{L}\)），从而 \(y^{\lambda} = f_{\beta}(x_{\beta}^{\lambda}) = 0\)（对所有 \(\lambda \in \mathbf{L}\)），这证明诸 \(\mathbf{M}^{\lambda}\) 之和是直和。

推论. 设 \((\mathbf{P}_{\alpha})\) 是 \(\mathbf{E}_{\alpha}\) 的子集的一个正向系统，并设 \(\mathbf{P} = \varinjlim \mathbf{P}_{\alpha}\)。若对所有 \(\alpha \in \mathbf{I}\)，\(\mathbf{P}_{\alpha}\) 是 \(\mathbf{E}_{\alpha}\) 的自由子集（相应地，基），则 \(\mathbf{P}\) 是 \(\mathbf{E}\) 的自由子集（相应地，基）。

第二个断言由第一个断言与命题4直接得出。因此只需证明：若诸 \(\mathbf{P}_{\alpha}\) 都是自由的，则每个子集 \(\{y^{(i)}\}_{1\leqslant i\leqslant n}\)（由 \(\mathbf{P}\) 中不同元素组成）都是自由的。存在一个 \(\alpha \in \mathbf{I}\) 和元素 \(x_{\alpha}^{(i)}\in \mathbf{P}_{\alpha}\)，使得 \(y^{(i)} = f_{\alpha}(x_{\alpha}^{(i)})\)（对 \(1\leqslant i\leqslant n\)）（集合论，III，§7，第5号，引理1）；若 \(\sum_{i}\lambda^{(i)}y^{(i)} = 0\)，则可设 \(\lambda^{(i)} = \phi_{\alpha}(\lambda_{\alpha}^{(i)})\)（对 \(1\leqslant i\leqslant n\)），于是 \(f_{\alpha}\left(\sum_{i}\lambda_{\beta}^{(i)}x_{\beta}^{(i)}\right) = 0\)；这蕴含 \(\sum_{i}\lambda_{\beta}^{(i)}x_{\beta}^{(i)} = 0\)（对某个 \(\beta \geqslant \alpha\)），其中 \(\lambda_{\beta}^{(i)} = \phi_{\beta \alpha}(\lambda_{\alpha}^{(i)})\)，\(x_{\beta}^{(i)} = f_{\beta \alpha}(x_{\alpha}^{(i)})\)，且诸 \(x_{\beta}^{(i)}\) 属于 \(\mathbf{P}_{\beta}\) 并两两不同，因为 \(y^{(i)} = f_{\beta}(x_{\beta}^{(i)})\)；于是 \(\lambda_{\beta}^{(i)} = 0\)（对 \(1\leqslant i\leqslant n\)），从而

\[
\lambda^ {(i)} = \phi_ {\beta} (\lambda_ {\beta} ^ {(i)}) = 0
\]

对 \(1 \leqslant i \leqslant n\)。

现在假设所有的环 \(\mathbf{A}_{\alpha}\) 都等于同一个环A，且 \(\phi_{\beta \alpha}\) 等于 \(1_{\mathbf{A}}\)；则对A-模的每个正向系统 \((\mathbf{E}_{\alpha}, f_{\beta \alpha})\)，\(\mathbf{E} = \varinjlim \mathbf{E}_{\alpha}\) 是一个A-模。设F是一个A-模，且对所有 \(\alpha\)，设 \(u_{\alpha}: \mathbf{E}_{\alpha} \to \mathbf{F}\) 是A-线性映射，使得 \((u_{\alpha})\) 是映射的一个正向系统；则 \(u = \varinjlim u_{\alpha}\) 是E到F的一个A-线性映射。反之，对每个A-线性映射 \(v: \varinjlim \mathbf{E}_{\alpha} \to \mathbf{F}\)，族 \(v_{\alpha} = v \circ f_{\alpha}\) 是A-线性映射的一个正向系统，使得 \(v = \varinjlim v_{\alpha}\)。另一方面，我们注意到对 \(\alpha \leqslant \beta\) 映射

\[
\operatorname{Hom} (f _ {\beta \alpha}, 1 _ {F}) = \bar {f} _ {\alpha \beta}: \operatorname{Hom} _ {A} (E _ {\beta}, F) \rightarrow \operatorname{Hom} _ {A} (E _ {\alpha}, F)
\]

是一个Z-模同态，使得 \((\mathrm{Hom}_{\mathbf{A}}(\mathrm{E}_{\alpha},\mathrm{F}),\bar{f}_{\alpha\beta})\) 是Z-模的一个逆向系统；由于 \(\bar{f}_{\alpha\beta}(v_{\beta}) = v_{\beta} \circ f_{\beta\alpha}\)，上述论述可表述如下：

命题6. 对由A-模组成的每个正向系统 \((\mathbf{E}_{\alpha}, f_{\beta \alpha})\) 以及每个A-模F，典范映射 \(u \mapsto (u \circ f_{\alpha})\) 是Z-模同构

\[
d _ {\mathrm{F}}: \operatorname{Hom} _ {\mathrm{A}} (\varinjlim \mathrm{E} _ {\alpha}, \mathrm{F}) \rightarrow \varprojlim \operatorname{Hom} _ {\mathrm{A}} (\mathrm{E} _ {\alpha}, \mathrm{F}).\tag{5}
\]

推论1. 对于每个A-模同态 \(v: \mathbf{F} \to \mathbf{F}'\) ，

\[
\bar {v} _ {\alpha} = \operatorname{Hom} \left(\mathrm{l} _ {\mathrm{E} _ {\alpha}}, v\right): \operatorname{Hom} \left(\mathrm{E} _ {\alpha}, \mathrm{F}\right)\rightarrow \operatorname{Hom} \left(\mathrm{E} _ {\alpha}, \mathrm{F} ^ {\prime}\right)
\]

构成Z-线性映射的一个逆向系统，且图表

\[
\begin{array}{c} \operatorname{Hom} (\varinjlim \mathrm{E} _ {\alpha}, \mathrm{F}) \xrightarrow {d _ {\mathrm{F}}} \varprojlim \operatorname{Hom} (\mathrm{E} _ {\alpha}, \mathrm{F}) \\ \operatorname{Hom} (1 _ {\mathrm{E}}, v) \Bigg \downarrow \qquad \qquad \qquad \qquad \qquad \qquad \qquad \qquad \qquad \qquad \qquad \qquad \qquad \qquad \qquad \varprojlim_ {\overline {{v}} _ {\alpha}} \\ \operatorname{Hom} (\varinjlim \mathrm{E} _ {\alpha}, \mathrm{F} ^ {\prime}) \xrightarrow {d _ {\mathrm{F} ^ {\prime}}} \varprojlim \operatorname{Hom} (\mathrm{E} _ {\alpha}, \mathrm{F} ^ {\prime}) \end{array}\tag{6}
\]

是交换的。

对所有 \(u \in \operatorname{Hom}(\varinjlim E_{\alpha}, F)\)，由定义有 \(d_{F'}(v \circ u) = (v \circ u \circ f_{\alpha})\)，而图(6)的交换性随后由定义直接得出。

推论2. 若 \((\mathbf{E}_{\alpha}, f_{\beta \alpha})\) 是左A-模的一个正向系统，\(\mathbf{E} = \varinjlim \mathbf{E}_{\alpha}\)，则 \((\mathbf{E}_{\alpha}^{*}, {}^t f_{\beta \alpha})\) 是右A-模的一个逆向系统，且 \(\varprojlim \mathbf{E}_{\alpha}^{*}\) 典范同构于 \(\mathbf{E}^{*}\)。

注记. 设E是一个A-模，且 \((\mathrm{M}_{\alpha})_{\alpha\in I}\) 是E的子模的一个递增族，使得E是诸 \(M_{\alpha}\) 之并；若 \(j_{\beta\alpha}:M_{\alpha}\to M_{\beta}\)（对 \(\alpha\leqslant\beta\)）且 \(j_{\alpha}:M_{\alpha}\to E\) 为典范单射，则显然 \(j=\varinjlim j_{\alpha}\) 是 \(\varinjlim M_{\alpha}\) 到E上的同构（集合论，III，§7，第 6 号，注记 1）。特别地，每个A-模都是其有限生成子模构成的右有向族的正向极限。

